\chapter{Conclusion and Limitations}
\label{sec:conclusion}

This thesis presents a hierarchical framework for Lean 4 theorem proving that separates local proof resolution from skeleton decomposition. The design addresses sparse reward through AST-guided recovery and handles decomposition credit through dependency-aware AND/OR backup. The key correctness boundary is that patched scripts with \texttt{sorry} are used only for learning signal; final success requires a fully verified Lean proof without placeholders.

The current design still has limitations. Dependency-based credit assignment reflects the proof that was found, not necessarily the shortest or mathematically minimal proof. Open subgoals within finite search limits cannot always be distinguished from genuinely impossible subgoals. The effectiveness of skeleton generation depends on the policy model's ability to propose useful subgoals. These limitations should be evaluated empirically in the experiments.

There are also broader limitations that are easier to state than to eliminate. The model interface assumes that the LLM can produce output that is at least parseable into Lean code blocks. If the model drifts into natural language, malformed markdown, or repetitive syntax errors, the Sorrifier can recover only a fraction of the damage. The scaffold-aware rule assumes that subgoals can be cleanly associated with a single target hole, which is true for the proof bodies GammaZero is designed to handle but may be awkward for very irregular generated code. The priority queue assumes that the heuristic score is at least informative enough to rank open states roughly by promise; if that score is poor, the system can still search, but it will search more expensively.

The dependency reward also has an important interpretive limit. It measures usage in the elaborated proof that was found, not a global theorem-theoretic notion of necessity. A declaration can be unused in the final stitched proof and still have been useful as an intermediate exploration step. Likewise, an unresolved child may block the final theorem even if it looks locally reasonable. These are not bugs in the reward definition; they are reminders that the thesis is about search data, not formal proof minimization.

Future work naturally follows from this framing. The collected search records could be used for policy fine-tuning or GRPO-style optimization. Retrieval could be integrated so that skeleton prompts and local proof prompts receive relevant premises automatically. Larger or more diverse theorem libraries could be added so that the search graph contains a broader variety of decomposition patterns. The skeleton scorer could also be replaced by a learned ranking model once enough trajectories have been collected. Each of these directions would build on the current thesis, but none of them is required for the present contribution: a verifiable, scaffold-aware, priority-queue proof-search system that produces structured reward data.

Another direction is to improve the model interface itself. The current separation between local proof and
skeleton prompts is simple and inspectable, and the implementation already feeds recent Lean failures
back into retry prompts. A future system could add richer search-history summaries, dependency
feedback, or retrieved examples of similar proof patterns. A model
might learn that certain goals should be attacked directly while others should be decomposed. It might
also learn to propose smaller skeletons that introduce only subgoals likely to be solved by the current
local proof policy. Such improvements would turn the current hand-designed local-proof-first rule into a learned
decision about when decomposition is useful.

The final long-term goal is a closed loop: generate search trajectories, score them with Lean-backed
signals, train a better proposal model, and then use the improved model to generate higher-quality
trajectories. This thesis implements the first half of that loop. It defines the graph structure, verification
requirements, repair mechanism, dependency analysis, and trajectory export needed for such training. The
remaining optimization step is left for future work because it requires substantial compute and careful
experimental design. Keeping that boundary explicit avoids overstating the contribution while still
showing how the system can support future reinforcement learning.

The thesis also leaves room for better proof-object analysis. The current dependency analyzer focuses on
whether generated declarations are used by the elaborated proof expression and whether they contain
placeholders. A richer analyzer could measure proof-term size, count references to library lemmas, detect
whether a declaration is used only through a trivial wrapper, or compare multiple stitched proofs for the
same theorem. Such analysis would help distinguish merely valid proofs from compact or pedagogically
useful proofs. That is outside the current scope, but the graph representation makes the extension
natural.

Finally, the system could benefit from richer feedback to the model at generation time. The current
implementation already feeds recent Lean failures back into retry prompts for the same state. Future
prompts could add higher-level summaries of reserved skeletons, dependency feedback from previous
decompositions, or retrieved examples from similar theorems. This would give the model more context
about the search history while preserving the same strict verification boundary: whatever the model
proposes, Lean must still check the final proof.
